\section{可检验命题与实验/数据策略}

\subsection{准周期谱与黄金支系判别}

\textbf{命题 P1（准周期谱线）。}若三维可观测集来自无理取向切投影模型集，则其衍射谱或结构因子应呈现离散纯点谱主导的准周期特征，并对余维相位偏移 $u$ 的变化表现为可预测的“相位滑移”而非随机漂移。\\
\textbf{命题 P2（黄金支系优选的统计信号）。}在固定资源约束下，若扫描斜率处于最小共振支系，则前缀读出序列的差分均匀性指标（如星差异或间隙分布）应在同复杂度类中达到可比较的极值表现，从而可用于与其他无理斜率进行统计区分。

\subsection{稳定扇区折叠的退化指纹}

\textbf{命题 P3（折叠退化直方图）。}在给定窗口长度 $m$ 与固定折叠规则下，$\Fold_m$ 的原像基数分布（退化直方图）为离散且可复现的组合不变量；不同实验平台只要共享相同的协议语法约束，其退化指纹应具有跨系统可迁移性。该命题可通过对离散读出记录进行窗口统计重构来检验。

\subsection{延迟选择的协议重构检验}

\textbf{命题 P4（等价类再分配）。}若延迟改变测量设置对应于改变读出核/投影参数，则观测统计的变化应可由“等价类分解的再分配”解释，并满足由折叠规则确定的可计算守恒关系（例如总权重守恒、特定类型间的可达迁移图约束）。该命题要求对实验读出进行类型层面的重编码，而非仅在波包图像层比较。

\subsection{自适应观测者的闭环稳定性}

\textbf{命题 P5（闭环收敛域）。}若观测者更新律近似执行失配泛函的下降，则存在可观察的收敛域与相变边界：当外部扰动或资源约束跨过阈值时，系统从稳定类型保持区进入缺陷累积区，并出现可测的失配增长与类型崩塌。可通过在可控系统（如量子光学反馈、可编程介质、或离散元胞自动机实验）中实施参数反馈并统计稳定扇区保持率检验。

